% LaTeX CV Template using Awesome CV
\documentclass[11pt, a4paper]{awesome-cv}

% Configure page margins
\geometry{left=1.5cm, top=1.0cm, right=1.5cm, bottom=1.5cm, footskip=1.0cm}

% Define custom colors
\definecolor{awesome}{HTML}{915DA3} % Custom dark blue for headers
\definecolor{highlight}{HTML}{1ABC9C} % Teal for highlights

% Header Information
\name{Alireza}{Amiri}
\position{Mechatronics Engineering Student}
\address{Tehran, Iran}
\mobile{+98 901 324 6001}
\email{alireza2001amiri@gmail.com}
\linkedin{linkedin.com/in/alirezaamiriii}
\quote{``Focused on tackling challenges with innovative solutions and boundless enthusiasm for learning."}

\begin{document}
	
	% Header
	\makecvheader[C]
	
	% Summary
	\cvsection{Objective}
Driven to solve challenges with creativity and determination, combining a passion for learning with the ability to navigate ambiguity and deliver results.
	
	\cvsection{Education}
	% Education
\begin{cventries}
	\cventry
	{M. S. Student in Mechatronics Engineering} % Degree
	{K. N. Toosi University of Technology University} % Institution
	{Tehran, Iran} % Location
	{Sep. 2023 -- Present} % Date(s)
	{
		\begin{cvitems} % Description(s) bullet points
			\item {GPA: 18.61/20}
			\item {Thesis: Design and Implementation of a Magnetically Levitated Planar Motor}
		\end{cvitems}
	}
	\cventry
	{Bachelor of Science in Electrical Engineering (Control Orientation)}
	{Shahed University}
	{Tehran, Iran}
	{Jul. 2019 -- Jul. 2023}
	{
		\begin{cvitems}
			\item {GPA: 16.03/20}
			\item {Thesis: Design and Implementation of a PID-based Control System for Achieving Stable Magnetic Levitation}
		\end{cvitems}
	}
	
	% Experience
	\cvsection{Experience}
	
		\cventry
	{Product Owner – Data Solutions and BI}
	{Moneyar}
	{Tehran, Iran}
	{Feb 2024 - Present}
	{
\begin{cvitems}
	\item {Led the research and development of data entry systems as foundational components for automated BI solutions.}
	\item {Designed and implemented processes for transforming raw data into actionable insights through advanced BI tools.}
	\item {Collaborated with cross-functional teams to explore innovative methodologies, ensuring cutting-edge data solutions.}
\end{cvitems}
	}
	
	
	\cventry
{Design And Development Assistant}
{Behnood Pharmed}
{Tehran, Iran}
{Aug 2023 - Oct 2024}
{
\begin{cvitems}
	\item {Gained practical insights into business processes by designing and developing management dashboards.}
	\item {Applied financial and project management concepts to real-world scenarios, enhancing professional growth.}
	\item {Utilized expertise in databases and QlikView to contribute to data-driven decision-making tools.}
\end{cvitems}
}

	\cventry
	{Maintenance and Repair Engineering Intern}
	{MAPNA Generator Engineering \& Manufacturing Co.}
	{Karaj, Iran}
	{Jun. 2022 -- Sep. 2022}
	{
		\begin{cvitems}
			\item {Gained insights into industrial large-scale factories and basic performance of industrial machinery.}
			\item {Learned fault diagnosis procedures and cost-efficient problem-solving techniques.}
			\item {Experimented with industrial control systems, including PLC and LOGO systems.}
		\end{cvitems}
	}
	
	\cventry
	{Translator}
	{Haftonim Publication}
	{Remote}
	{Mar. 2021 -- Jul. 2021}
	{
		\begin{cvitems}
			\item {Translated a scientific book, "Archimedes and the Door of Science," for teenage readers.}
			\item {Developed skills in literary translation and adapted text to Persian.}
		\end{cvitems}
	}
	
	\cventry
	{NFT Marketing}
	{Self-Employed}
	{Remote}
	{Jan. 2022 -- Jul. 2022}
	{
		\begin{cvitems}
			\item {Launched and marketed "Callipen," a calligraphy artwork brand.}
			\item {Achieved international sales and developed cross-cultural communication skills.}
		\end{cvitems}
	}
	
	% Projects
		\cvsection{Selected Projects}
		
	\cventry
	{A Project of Mechatronics course}
	{Waypoint Generation based on Crop Row Detection Using Unet and Hough Transform}
	{}
	{Jan. 2024 -- Sep. 2024}
	{
		\begin{cvitems}
			\item {Develops a method to generate precise waypoints for autonomous agricultural robots using aerial imagery.}
			\item {Utilizes U-Net based image segmentation for crop detection and Hough Transform for row detection, followed by a mathematical conversion to global coordinates.}
			\item {The proposed system achieved high accuracy in crop and row detection, enabling precise waypoint generation for autonomous robot navigation.}
		\end{cvitems}
	}

	\cventry
	{B.Sc. Thesis}
	{Magnetic Levitation}
	{}
	{Mar. 2022 -- Sep. 2023}
	{
		\begin{cvitems}
			\item {Developed an Arduino-based PID controller to achieve stable magnetic levitation.}
			\item {Experimented with PID and H-infinity control methods using MATLAB and Simulink.}
			\item {Designed and optimized an electrical circuit to support system stability.}
		\end{cvitems}
	}
	\cventry
	{Digital Signal Processing Project}	
	{Speech Emotion Detection}
	{}
	{Jul. 2022 -- Aug. 2022}
	{
		\begin{cvitems}
			\item {Trained a machine learning model to classify emotions and gender from Persian speech data.}
			\item {Achieved 60\% accuracy in the best-case scenario.}
		\end{cvitems}
	}
	
%	\cventry
%	{ECG Signal Monitoring}
%	{Biomedical Project}
%	{}
%	{Jul. 2023 -- Jul. 2023}
%	{
%		\begin{cvitems}
%			\item {Designed an ESP32-based system for recording ECG signals and heart rates.}
%			\item {Collaborated on biomedical device development, learning about medical sensors.}
%		\end{cvitems}
%	}
\cvsection{Publications}	
\begin{cvskills}
	\cvskill
	{Waypoint Generation based on Crop Row Detection Using Unet and Hough Transform}
	{Accepted paper in ICRoM 2024}
	% Skills
	\end{cvskills}
	
\cvsection{Skills}
\begin{cvskills}
\cvskill
{Programming Languages}
{Python,Arduino, Verilog, SQL, C++}
\cvskill{Engineering Tools}{MATLAB, Simulink, Proteus, Ansys, ISE}
\cvskill{Soft Skills}{Creativity, Problem Solving, System Thinking,  Adaptibily, Teamwork}
\end{cvskills}	
	% Languages
\cvsection{Languages}

\begin{cvskills}
\cvskill{Persian}{Native}
\cvskill{English}{Fluent}
\cvskill{Arabic}{Conversational}
\end{cvskills}	
	% Footer
	\makecvfooter
	{\today}
	{Alireza Amiri~~~·~~~CV}
	{\thepage}
\end{cventries}	
\end{document}
